\scanpage{56}
condition nécessaire évidente pour l'existence d'une rev. quadr. $\Delta$ de $S$ correspondant à la donnée $(L,D)$], et on trouve alors, en \deleted{\unsure{fct|fonctions}} \inserted{termes} \inserted{de $(L,b_0)$, le couple} $(F,\underline b\in\Gamma(S,\underset{\scriptscriptstyle\Vert\,2}{f}F))$, donc $(B,e_B,\mathrm{tr})$.

Les structures \deleted{d'algèbre} cherchées \inserted{$q$} sur $B$ correspondent donc aux sections $\underline c$ d'un \inserted{certain} torseur $Q(\underline b)$ sous $\check{\mathbb V}(L^{\otimes2})$. \deleted{Mais \unsure{on}} La fonction \deleted{$\unsure{\Theta(}$} $\Theta:q\longmapsto\deleted{\unsure{\Theta(q)}}=D(q)=-\mathrm{discr}(q)$ \inserted{de $Q(\underline b)$ dans $\check{\mathbb V}(L^{\otimes2})$} \deleted{correspondant donc aux} \inserted{est alors une} fonction linéaire affine, qui satisfait
\[\deleted{\unsure{D}}\Theta(q+\underset{\substack{\uparrow\\\Gamma(L^{\otimes2})}}{\varpi})=\Theta(q)-4\varpi\tag*{(107)}\]
% slanted
\margin{\uline{\unsure{remonter|renvois|remarques}} \uline{\unsure{au|aux} § 6}}
i.e. la partie homogène $\Theta_0$ de $\Theta$ est définie par
\[\Theta_0(\varpi)=-4\varpi\qquad\text{i.e.}\quad\Theta_0=-4\,\mathrm{id}\tag*{(108)}\]

La section $q$ cherchée, doit donc satisfaire la condition
\[\Theta(q)=D\tag*{(109)}\]
ce qui, en termes d'une trivialisation locale $q_0$ de $Q(\underline b)$, permettant d'écrire $q$ comme
\[q=q_0+\varpi\qquad\varpi\in\Gamma(S,L^{\otimes2})\tag*{(110)}\]
s'écrit comme
\[\Theta(q_0)\unsure{-|+|\dot{-}}4\varpi=D\qquad\text{i.e.}\quad4\varpi=\Theta(q_0)-D)\tag*{(111)}\]
% archive sheet number: 55
